% Section 5.9, from lectures/L05/appendix/b_exercises.md.

\section{Uppgifter}
\label{c5:sec:uppgifter}

\begin{aside}
\textbf{Så arbetar du med uppgifterna.} Del~1 och del~2 räknas under lektionen: först på egen
hand, några minuter per uppgift, sedan gemensamt. Del~3 är extrauppgifter att träna vidare på
efter lektionen. De flesta går att räkna i huvudet eller med papper och penna.

\facitnotis{L05}
\end{aside}

% ------------------------------------------------------------------------------------------
\uppgiftsdel{Del 1}{Repetitionsuppgifter}

\uppgift{1.1}{Ekvationssystem: spänning och ström}
Lös följande ekvationssystem (strömmar i ampere):
\[
  \begin{cases}
    2I_1 + 3I_2 = 13 \\
    4I_1 - I_2 = 5
  \end{cases}
\]

\uppgift{1.2}{Kretsanalys med KVL och KCL}
En krets har tre resistorer i ett nätverk. KVL och KCL ger följande system för $U_A$ och $U_B$
(nodspänningar i volt):
\[
  \begin{cases}
    3U_A - U_B = 12 \\
    -U_A + 4U_B = 6
  \end{cases}
\]
Lös systemet och ange $U_A$ och $U_B$.

% ------------------------------------------------------------------------------------------
\uppgiftsdel{Del 2}{Nytt stoff}

\uppgift{2.1}{Potensregler}
Förenkla med hjälp av potensreglerna. Lämna svaret med positiva exponenter.
\begin{delar}(5)
  \task $x^3 \cdot x^5$
  \task $\dfrac{R^6}{R^2}$
  \task $(I^2)^3$
  \task $U^{-2} \cdot U^5$
  \task $\left(\dfrac{2}{R}\right)^3$
\end{delar}

\uppgift{2.2}{Standardform för komponentvärden}
Skriv följande värden på standardform, $a \times 10^n$ med $1 \leq a < 10$.
\begin{delar}(4)
  \task $47\,000\,\Omega$
  \task $0{,}000\,022\,\text{F}$
  \task $3\,300\,000\,\text{Hz}$
  \task $0{,}000\,000\,015\,\text{A}$
\end{delar}

\uppgift{2.3}{Effektberäkning}
Beräkna effekten $P$ som utvecklas i ett motstånd med hjälp av $P = \dfrac{U^2}{R}$ och
$P = RI^2$.
\begin{parts}
  \item $U = 5\,\text{V}$ och $R = 100\,\Omega$. Beräkna $P$ i mW.
  \item $I = 40\,\text{mA}$ och $R = 50\,\Omega$. Beräkna $P$ i mW.
\end{parts}

\uppgift{2.4}{RMS-värde för sinusspänning}
Sambandet mellan amplituden $\abs{U}$ och effektivvärdet (RMS-värdet) $U_{\text{RMS}}$ för en
sinusspänning är
\[
  U_{\text{RMS}} = \frac{\abs{U}}{\sqrt{2}}.
\]
\begin{parts}
  \item En sinusspänning har amplituden $\abs{U} = 325\,\text{V}$ (nätspänning). Beräkna
    $U_{\text{RMS}}$.
  \item En uppmätt RMS-spänning är $U_{\text{RMS}} = 12\,\text{V}$. Beräkna amplituden $\abs{U}$.
\end{parts}

% ------------------------------------------------------------------------------------------
\uppgiftsdel{Del 3}{Extrauppgifter}
Uppgifterna nedan är extra träning och görs med fördel på egen hand efter lektionen.

\uppgift{3.1}{Beräkna potenser}
Beräkna \textbf{utan} miniräknare.
\begin{delar}(6)
  \task $3^4$
  \task $(-2)^5$
  \task $10^{-2}$
  \task $5^0$
  \task $2^{-3}$
  \task $0{,}1^3$
\end{delar}

\uppgift{3.2}{Potensregler}
Förenkla. Svaret ska ha positiva exponenter.
\begin{delar}(6)
  \task $a^7 \cdot a^{-3}$
  \task $\dfrac{U^5}{U^8}$
  \task $(2R^3)^2$
  \task $\dfrac{(I^2)^4}{I^3}$
  \task $\left(\dfrac{R^2}{3}\right)^3$
  \task $\dfrac{6x^5}{2x^2}$
\end{delar}

\uppgift{3.3}{Tiopotenser}
Skriv som en enda tiopotens.
\begin{delar}(5)
  \task $10^4 \cdot 10^{-7}$
  \task $\dfrac{10^{-2}}{10^{-5}}$
  \task $\left(10^3\right)^{-2}$
  \task $\dfrac{10^6 \cdot 10^{-2}}{10^{-3}}$
  \task $\sqrt{10^{-6}}$
\end{delar}

\uppgift{3.4}{Rötter}
\begin{parts}
  \item Beräkna $\sqrt{81}$.
  \item Beräkna $\sqrt{0{,}25}$.
  \item Beräkna $\sqrt[3]{27}$.
  \item Beräkna $\sqrt{16 \times 10^{-4}}$.
  \item Förenkla $\sqrt{50}$ till exakt form.
  \item Visa med $a = 9$ och $b = 16$ att $\sqrt{a + b} \neq \sqrt{a} + \sqrt{b}$.
\end{parts}

\uppgift{3.5}{Bråkexponenter}
Beräkna.
\begin{delar}(5)
  \task $16^{1/2}$
  \task $27^{1/3}$
  \task $9^{3/2}$
  \task $8^{-1/3}$
  \task $\left(\dfrac{1}{4}\right)^{-1/2}$
\end{delar}

\uppgift{3.6}{Från prefix till grundenhet}
Skriv om till grundenheten på standardform.
\begin{delar}(6)
  \task $2{,}2\,\text{k}\Omega$
  \task $470\,\mu\text{F}$
  \task $15\,\text{pF}$
  \task $250\,\text{mA}$
  \task $1{,}8\,\text{MHz}$
  \task $6{,}8\,\text{nH}$
\end{delar}

\uppgift{3.7}{Från standardform till prefix}
Ange värdet med det prefix som ger ett tal mellan $1$ och $1000$.
\begin{delar}(3)
  \task $3{,}3 \times 10^{-6}\,\text{F}$
  \task $4{,}7 \times 10^{5}\,\Omega$
  \task $2{,}5 \times 10^{-4}\,\text{A}$
  \task $8 \times 10^{7}\,\text{Hz}$
  \task $1{,}2 \times 10^{-11}\,\text{F}$
\end{delar}

\uppgift{3.8}{Räkning med standardform}
Beräkna och ange svaret på standardform.
\begin{delar}(3)
  \task $\left(4 \times 10^5\right)\left(2 \times 10^{-8}\right)$
  \task $\dfrac{9 \times 10^{-3}}{3 \times 10^{2}}$
  \task $\left(5 \times 10^3\right)^2$
  \task $2{,}5 \times 10^{-3} + 7{,}5 \times 10^{-4}$
  \task $\dfrac{\left(6 \times 10^{-2}\right)^2}{4 \times 10^{-5}}$
\end{delar}

\uppgift{3.9}{Värdesiffror och avrundning}
\begin{parts}
  \item Hur många värdesiffror har $4{,}70 \times 10^3$?
  \item Avrunda $0{,}024\,68$ till två värdesiffror.
  \item Avrunda $12\,345\,\Omega$ till tre värdesiffror och skriv svaret på standardform.
  \item En spänning mäts till $U = 5{,}0\,\text{V}$ och en resistans till $R = 100\,\Omega$ (tre
    värdesiffror). Med hur många värdesiffror bör strömmen $I = \dfrac{U}{R}$ anges? Beräkna
    strömmen.
\end{parts}

\uppgift{3.10}{Ohms lag med tiopotenser}
Beräkna och ange svaret med lämpligt prefix.
\begin{parts}
  \item $U = 3{,}3\,\text{V}$ och $R = 220\,\Omega$. Beräkna $I$.
  \item $I = 250\,\mu\text{A}$ och $R = 47\,\text{k}\Omega$. Beräkna $U$.
  \item $U = 12\,\text{V}$ och $I = 4\,\text{mA}$. Beräkna $R$.
\end{parts}

\uppgift{3.11}{Effekt med tiopotenser}
\begin{parts}
  \item $R = 1{,}5\,\text{k}\Omega$ och $I = 20\,\text{mA}$. Beräkna $P = RI^2$.
  \item $U = 5\,\text{V}$ och $R = 2{,}2\,\text{k}\Omega$. Beräkna $P = \dfrac{U^2}{R}$ i mW.
  \item En resistor är märkt $470\,\Omega$ och $0{,}25\,\text{W}$. Beräkna den största tillåtna
    strömmen.
  \item Beräkna den största tillåtna spänningen över resistorn i c).
\end{parts}

\uppgift{3.12}{Tidskonstant med prefix}
Tidskonstanten för en RC-krets är $\tau = RC$.
\begin{parts}
  \item $R = 10\,\text{k}\Omega$ och $C = 4{,}7\,\mu\text{F}$. Beräkna $\tau$ i ms.
  \item $\tau = 1\,\text{ms}$ och $C = 100\,\text{nF}$. Beräkna $R$.
  \item $R = 2{,}2\,\text{M}\Omega$ och $\tau = 1{,}1\,\text{s}$. Beräkna $C$ och ange svaret i nF.
\end{parts}

\uppgift{3.13}{Toppvärde, effektivvärde och effekt}
För en sinusspänning gäller $U_{\text{RMS}} = \dfrac{\abs{U}}{\sqrt{2}}$.
\begin{parts}
  \item Beräkna $U_{\text{RMS}}$ då amplituden är $\abs{U} = 10\,\text{V}$.
  \item Beräkna amplituden $\abs{U}$ då $U_{\text{RMS}} = 230\,\text{V}$.
  \item Effekten i ett motstånd beräknas ur effektivvärdet, $P = \dfrac{U_{\text{RMS}}^2}{R}$.
    Visa att detta kan skrivas $P = \dfrac{\abs{U}^2}{2R}$.
  \item Beräkna $P$ då $\abs{U} = 12\,\text{V}$ och $R = 8\,\Omega$.
\end{parts}

\uppgift{3.14}{Hitta felet}
Varje rad innehåller ett vanligt potensfel. Förklara felet och skriv det korrekta svaret.
\begin{delar}(3)
  \task $2^3 \cdot 2^4 = 2^{12}$
  \task $\left(a^2\right)^3 = a^5$
  \task $\dfrac{10^6}{10^2} = 10^3$
  \task $(2a)^3 = 2a^3$
  \task $\sqrt{9 + 16} = 3 + 4$
  \task $a^{-2} = -a^2$
\end{delar}
